\section{Singular acquired laws and changing strata}\label{sec:strata}
The density assumption may be replaced by an explicitly stratified acquired law. Allow normalized nonnegative report kernels, including perfectly informative marked reports. Bayes updates are defined on the events of positive report probability; arbitrary definitions on null events have no effect. Let $\mathcal S$ be a finite collection of distinct simplex faces, including possible vertices. For $s\in\mathcal S$ write $d_s$ for its affine dimension. Fix a compact convex $K_s$ in the relative interior of that face, with a fixed relative-interior margin. Different relative interiors are disjoint, so the stratum is a Borel function of the posterior. It is not supplied as an uncharged retained register.

\begin{theorem}[Stratified posterior pooling]\label{thm:strata}
Suppose uniformly in $p,a$ that every actual one-step posterior law decomposes as
\begin{equation}\label{eq:stratumlaw}
 \Lambda_{t-1}(p,a)=\sum_{s\in\mathcal S}\Lambda_{t,s}(p,a),
 \qquad \operatorname{supp}\Lambda_{t,s}(p,a)\subset K_s.
\end{equation}
For $d_s>0$ suppose $\Lambda_{t,s}(p,a)$ has relative-coordinate density at most $\rho_{t,s}$. These are subprobability laws: $\rho_{t,s}$ includes their actual mass. A zero-dimensional component is supported on its one vertex.

Choose sufficiently large integers $L_{t,s}$ for each positive-dimensional stratum and one label for each vertex. There is a causal observer of risk
\begin{equation}\label{eq:stratifiedupper}
 R-B_n^*\le C\sum_{t=1}^n\sum_{s:d_s>0}\rho_{t,s}L_{t,s}^{-2/d_s}
\end{equation}
with at most
\begin{equation}\label{eq:stratumstates}
 1+\sum_{t=1}^n\left(\#\{s:d_s=0\}+\sum_{s:d_s>0}L_{t,s}\right)
\end{equation}
total states. Its excess equals the stratumwise version of \eqref{eq:poolidentity}. The constants depend only on the finite relative geometries and the terminal loss.

If a terminal stratum $s_*$ of dimension $r>0$ has mass at least $w>0$ under \emph{every} conditional last-step law and density at most $\rho_*$, then every checkpoint or autonomous $M$-state procedure has
\begin{equation}\label{eq:stratumlower}
 R-B_n^*\ge c_r w^{1+2/r}\rho_*^{-2/r}M^{-2/r}.
\end{equation}
In particular, if all other stratum dimensions are at most $r$, checkpoint and autonomous memory exponents match at $2/r$ for each fixed horizon. No uniformity is claimed as $w$ vanishes or a relative-interior margin degenerates.
\end{theorem}
\begin{proof}
Keep distinct labels for different faces, and partition each positive-dimensional $K_s$ by a grid in its affine coordinates. Pool all incoming laws in each of these cells. Their conditional barycenters stay in the same compact convex face subset. Lemma~\ref{lem:aggregation} therefore gives exact retained posteriors just as before. Restrict the concave $V_t$ to that face; it remains uniformly Lipschitz in relative coordinates. Lemma~\ref{lem:integrated} applies to each actual submeasure without renormalization. Vertex components have zero Jensen defect because all their posterior values are the same vertex. Sum the defects to obtain \eqref{eq:stratifiedupper}. The labels $(t,s,j)$ and the initial state give exactly the count \eqref{eq:stratumstates}, so the phase and stratum are not free.

Under any full-history policy the terminal $s_*$-submeasure still has mass at least $w$ and density at most $\rho_*$. Project any $M$ centers onto its affine span. Their radius-$u$ balls cover mass at most $\rho_*v_rMu^r$. Choosing $u=(w/(2\rho_*v_rM))^{1/r}$ leaves at least $w/2$ mass uncovered. Projection and conditional squared-loss decomposition give \eqref{eq:stratumlower}. An equal constant-fraction allocation of the available labels to the finitely many stages and positive-dimensional strata gives exponent $2/r$ in the upper, since $L^{-2/d_s}\le L^{-2/r}$ for $d_s\le r$. The fixed-horizon checkpoint upper follows by applying the same stratified terminal grid to an optimal full-history policy. This proves both matching assertions without deriving an online lower from terminal geometry when an intermediate dimension is larger than $r$.
\end{proof}

\begin{corollary}[Intermittently validated controlled sensing]\label{cor:validation}
At call $t$ of the two-state sensor in Theorem~\ref{thm:sensors}, independently choose a validated report with calibrated probability $\theta_t\in[0,1]$. A validated report physically reveals the new hidden state; otherwise the selected overlapping sensor emits its continuous report. The validation flag is part of this single paid report. Then the acquired posterior law is the sum of vertex atoms of total mass $\theta_t$ and a continuous component of mass exactly $1-\theta_t$ and density at most $(1-\theta_t)\rho$.

A controller using $L_t$ continuous cells and two vertex labels per stage has $1+\sum_t(L_t+2)$ total states and
\begin{equation}\label{eq:validationupper}
 R-B_n^*\le C\rho\sum_{t=1}^n(1-\theta_t)L_t^{-2}.
\end{equation}
If $\theta_n<1$, both checkpoint and autonomous excesses are bounded below by
\begin{equation}\label{eq:validationlower}
 c(1-\theta_n)\rho^{-2}M^{-2}.
\end{equation}
For fixed $n$ with $\theta_n$ bounded away from one, the memory exponent is two. No lower floor survives when the terminal call validates the state with probability one.
\end{corollary}
\begin{proof}
On a validated report the posterior is exactly $0$ or $1$, not a formal limiting state. On a nonvalidated report its law is that of Theorem~\ref{thm:sensors}. The independent validation probability gives the stated actual masses. A positive physical transition carries a vertex posterior to a predicted probability in $[\alpha,1-\alpha]$, so the next noisy report can again acquire a continuous posterior. Thus the experiment repeatedly moves between zero- and one-dimensional acquired strata. Apply Theorem~\ref{thm:strata} with two vertices and the compact noisy interval. In \eqref{eq:stratumlower}, substitute $w=1-\theta_n$ and $\rho_*=(1-\theta_n)\rho$ to obtain \eqref{eq:validationlower}.
\end{proof}
For two calls, if both validation probabilities are less than one, the strict feedback gain from Theorem~\ref{thm:sensors} remains at least $(1-\theta_1)(1-\theta_2)$ times its displayed bound. The nonvalidated first-report event already gives that gain; validated histories give a nonnegative additional benefit to the optimized feedback policy. This singular extension does not leave changing rank as an interface assertion, and rare-stratum mass is visible in both bounds.
